\begin{abstract}
Este trabalho apresenta a implementação e avaliação experimental de um sistema
distribuído que executa o algoritmo de eleição Bully e utiliza heartbeats
periódicos para detectar falhas do coordenador. Os nós são executados como
processos independentes em contêineres Docker e comunicam-se exclusivamente via
gRPC com deadlines explícitos. Além do cenário base com três nós e variação do
intervalo de heartbeat, a solução foi avaliada em uma bateria modular com clusters
de $N \in \{3, 5, 7, 10\}$ nós, investigando a complexidade de mensagens
$\mathcal{O}(N^2)$, a tolerância a falhas simultâneas dos maiores nós e a estabilidade
diante de falhas intermediárias. Os resultados confirmaram 100\% de sucesso em
todos os cenários, demonstrando a constância da latência de convergência
($\approx 1{,}15\,\text{s}$) e o crescimento quadrático na troca de mensagens de controle.
\end{abstract}

\begin{IEEEkeywords}
Sistemas distribuídos, algoritmo Bully, eleição de líder, detector de falhas,
gRPC, Docker Compose, complexidade de mensagens.
\end{IEEEkeywords}
